С развитием цифровых технологий и повсеместным распространением мобильных устройств
взаимодействие горожан с городской средой приобретает новые формы.
Жители, ежедневно перемещающиеся по городу, первыми замечают изменения инфраструктуры,
аварийные ситуации и иные события, требующие оперативного реагирования со стороны
муниципальных служб. В то же время оперативная фиксация подобных событий и передача
информации ответственным органам остаётся затруднённой: гражданину необходимо
самостоятельно определить ответственную службу, найти её контактные данные и составить
обращение по телефону или электронной почте. Единый цифровой канал для подачи обращений,
как правило, отсутствует, а обратная связь о статусе обращения не предоставляется.

Решение данной проблемы требует комплексного подхода, охватывающего весь цикл
обработки обращения~--- от момента фиксации проблемы жителем на месте до контроля
её устранения диспетчером. Такая система должна объединять мобильное приложение,
позволяющее жителю оперативно зафиксировать проблему с фотографией и геолокацией,
серверную инфраструктуру для хранения и обработки данных и веб-панель диспетчера
для централизованного управления обращениями на интерактивной карте города.

Актуальность подтверждается следующими факторами:
\begin{itemize}
    \item растущие требования к оперативности сбора и обработки обращений граждан
          по вопросам городской инфраструктуры;
    \item необходимость обеспечения прозрачности процесса обработки обращений и
          информирования заявителей о текущем статусе;
    \item потребность муниципальных служб в централизованном приёме и диспетчеризации
          обращений, поступающих от граждан через мобильные приложения;
    \item повсеместное распространение мобильных устройств, позволяющих фиксировать
          проблемы непосредственно на месте их обнаружения;
    \item отсутствие единого решения, объединяющего фиксацию проблемы, диспетчеризацию
          на карте и маршрутизацию к ответственной службе, применимого к городу Кирову.
\end{itemize}

Целью работы является разработка программной системы фиксации городских событий,
обеспечивающей полный цикл обработки обращений~--- от подачи жителем через мобильное
приложение до контроля исполнения диспетчером через веб-панель. Система ориентирована
на централизованную обработку обращений жителей о проблемах городской инфраструктуры
города Кирова.

Для достижения данной цели необходимо решить следующие задачи:
\begin{enumerate}
    \item провести анализ предметной области и сформулировать требования к системе;
    \item выполнить обзор аналогов~--- существующих систем для фиксации проблем
          городской инфраструктуры;
    \item спроектировать структуру программной системы, включающей серверную часть,
          диспетчерскую веб-панель и мобильное приложение для жителей;
    \item спроектировать структуру базы данных и пользовательские интерфейсы
          веб-панели и мобильного приложения;
    \item реализовать серверную часть с программным интерфейсом, базой данных
          и алгоритмами дедупликации и маршрутизации обращений;
    \item реализовать клиентские приложения~--- диспетчерскую веб-панель с интерактивной
          картой и кроссплатформенное мобильное приложение для жителей;
    \item провести тестирование всех компонентов программной системы.
\end{enumerate}
